% Preamble for "Tamil, Step by Step, Root by Root" (HL00 book format).
% XeLaTeX only. Uses the vendored static Noto Sans Tamil (see _fonts/),
% loaded by relative path so local and CI builds are identical.
% (Non-polyglossia packages first, as a habit, matching the other tracks.)

\usepackage{fontspec}
\setmainfont{Latin Modern Roman}
\setsansfont{Latin Modern Sans}
\setmonofont{Latin Modern Mono}

% Two Tamil-transliteration letters (alveolar ṉ and ṟ, "n/r with line below")
% are absent from Latin Modern. Render them with the bar-under accent, which
% works in any font, so the romanizations don't drop characters.
\usepackage{newunicodechar}
\newunicodechar{ṉ}{\b{n}}
\newunicodechar{ṟ}{\b{r}}
\newunicodechar{ḻ}{\b{l}}
\newunicodechar{ṁ}{\.{m}}
\newunicodechar{ḱ}{\'{k}}

\usepackage[table]{xcolor}
\usepackage{tcolorbox}
\tcbuselibrary{skins,breakable}
\usepackage{booktabs}
\usepackage{longtable}
\usepackage{tabularx}
\usepackage{array}
\usepackage[hidelinks]{hyperref}
\usepackage{titlesec}

\usepackage{polyglossia}
\setmainlanguage{english}
\setotherlanguage{tamil}
\setotherlanguage{arabic}

% Vendored Tamil font (../../_fonts/ from <lang>/book/).
\newfontfamily\tamilfont[Path=../../_fonts/, Script=Tamil]{NotoSansTamil-Static.ttf}

% Inline Tamil snippet within English text.
\newcommand{\ta}[1]{\texttamil{#1}}
\newfontfamily\telugufont[Path=../../_fonts/, Script=Telugu]{NotoSansTelugu-Static.ttf}
\newcommand{\te}[1]{{\telugufont #1}}
\newfontfamily\kannadafont[Path=../../_fonts/, Script=Kannada]{NotoSansKannada-Static.ttf}
\newcommand{\ka}[1]{{\kannadafont #1}}
\newfontfamily\malayalamfont[Path=../../_fonts/, Script=Malayalam]{NotoSansMalayalam-Static.ttf}
\newcommand{\ml}[1]{{\malayalamfont #1}}
\newfontfamily\devanagarifont[Path=../../_fonts/, Script=Devanagari]{NotoSansDevanagari-Static.ttf}
\newcommand{\dv}[1]{{\devanagarifont #1}}
\newfontfamily\arabicfont[Path=../../_fonts/, Script=Arabic]{NotoNaskhArabic-Static.ttf}
\newcommand{\ar}[1]{\textarabic{#1}}

\hypersetup{
  pdftitle={Tamil, Step by Step, Root by Root},
  pdfauthor={Adhithya Rajasekaran},
  pdfsubject={Etymology-driven Tamil curriculum},
  pdfkeywords={Tamil, Dravidian, language learning, etymology}
}

\newtcolorbox{cousinweb}{
  breakable, skin=enhanced, colback=yellow!8, colframe=yellow!45!black,
  boxrule=0.5pt, arc=1mm, left=6pt, right=6pt, top=4pt, bottom=4pt,
  fonttitle=\bfseries, title={The word, taken apart --- its roots and family}
}
\newtcolorbox{culture}{
  breakable, skin=enhanced, colback=red!5, colframe=red!40!black,
  boxrule=0.5pt, arc=1mm, left=6pt, right=6pt, top=4pt, bottom=4pt,
  fonttitle=\bfseries, title={Why it's said this way}
}
\newtcolorbox{grammarlens}[1][]{
  breakable, skin=enhanced, colback=blue!5, colframe=blue!35!black,
  boxrule=0.5pt, arc=1mm, left=6pt, right=6pt, top=4pt, bottom=4pt,
  fonttitle=\bfseries, title={Grammar Lens}, #1
}
\newtcolorbox{sounds}{
  breakable, skin=enhanced, colback=green!6, colframe=green!30!black,
  boxrule=0.5pt, arc=1mm, left=6pt, right=6pt, top=4pt, bottom=4pt,
  fonttitle=\bfseries, title={The letters in this word}
}
% Dravidian cognate box: the same idea across the family, forms supplied in full.
\newtcolorbox{cognates}{
  breakable, skin=enhanced, colback=violet!6, colframe=violet!45!black,
  boxrule=0.5pt, arc=1mm, left=6pt, right=6pt, top=4pt, bottom=4pt,
  fonttitle=\bfseries, title={Across the family --- the same idea, five ways}
}

\titleformat{\chapter}[display]
  {\normalfont\Large\bfseries}{Chapter \thechapter}{10pt}{\huge}
